\documentclass[11pt]{article}
\usepackage[margin = 2.0cm]{geometry}

\usepackage[utf8]{inputenc}
\usepackage{amsmath,amssymb,bm}

\begin{document}

\section*{Hamiltoniano efectivo para una cadena SSH acoplada por un fotón}

Partimos del hamiltoniano

\begin{equation}
H = H_{SSH}(t_1, t_2) + g(a + a^\dagger) H_{SSH}(t_1, -t_2)
\end{equation}

con condiciones periódicas. Diagonalizamos la parte puramente material:

\begin{equation}
H_{SSH}(t_1, t_2)
=
\sum_k |\Delta_0(k)| \,
\psi_k^\dagger \hat{\sigma}_z^k \psi_k
=
\sum_{k,\sigma}
\epsilon_{k,\sigma}
\hat{c}_{k,\sigma}^\dagger
\hat{c}_{k,\sigma}
\end{equation}

con

\begin{equation}
\Delta_0(k) = t_1 + t_2 e^{-ik},
\end{equation}

\begin{equation}
\epsilon_{k,\sigma} = (-1)^\sigma |\Delta_0(k)|.
\end{equation}

Para esta diagonalización se hicieron dos pasos. Primero un cambio de gauge:

\begin{equation}
\hat{c}_{k,A}^\dagger \hat{c}_{k,B}
\rightarrow
e^{-i\varphi_0(k)}
\hat{c}_{k,A}^\dagger \hat{c}_{k,B}
\end{equation}

con

\begin{equation}
\varphi_0(k)
=
\operatorname{Arg}\left(\Delta_0(k)\right).
\end{equation}

Segundo, un cambio de variables:

\begin{equation}
\hat{c}_{k,A}
=
\frac{
\hat{d}_{k,A} + \hat{d}_{k,B}
}{\sqrt{2}}
\end{equation}

\begin{equation}
\hat{c}_{k,B}
=
\frac{
\hat{d}_{k,A} - \hat{d}_{k,B}
}{\sqrt{2}}
\end{equation}

En este cambio de variables las matrices de Pauli transforman a:

\begin{equation}
\sigma_x^k \rightarrow \sigma_z^k
\end{equation}

\begin{equation}
\sigma_y^k \rightarrow -\sigma_y^k
\end{equation}

\begin{equation}
\sigma_z^k \rightarrow \sigma_x^k
\end{equation}

Tras todo esto, el hamiltoniano de acoplamiento transformó a:

\begin{equation}
H_{SSH}(t_1,-t_2)
=
\sum_k
\psi_k^\dagger
\left(
\Delta_z(k)\hat{\sigma}_z^k
+
\Delta_y(k)\hat{\sigma}_y^k
\right)
\psi_k
\end{equation}

con

\begin{equation}
\Delta_z(k)
=
\frac{
t_1^2 - t_2^2
}{
|\Delta_0(k)|
}
\end{equation}

\begin{equation}
\Delta_y(k)
=
\frac{
2t_1 t_2 \sin(k)
}{
|\Delta_0(k)|
}.
\end{equation}

Se puede mostrar que el módulo de $\Delta_y(k)$ se anula en
$k = 0,\pi$, es una función impar y además tiene un máximo para
$\pi/2 < |k| < \pi$.

A continuación realizamos una transformación de Schrieffer--Wolff
para eliminar la parte diagonal en materia de este acoplamiento.
Definimos:

\begin{equation}
\hat{H} = \hat{H}_0 + \hat{V}
\end{equation}

con

\begin{equation}
\hat{V}
=
\sum_{k,\sigma}
\Delta_z^\sigma(k)
\hat{c}_{k,\sigma}^\dagger
\hat{c}_{k,\sigma},
\end{equation}

\begin{equation}
\Delta_z^\sigma(k)
=
(-1)^\sigma \Delta_z(k),
\end{equation}

y la parte que queremos conservar del hamiltoniano:

\begin{equation}
\hat{H}_0
=
H_{SSH}(t_1,t_2)
+
\hbar\omega_c \hat{a}^\dagger \hat{a}
+
g(a+a^\dagger)
\underbrace{
\sum_k
i\Delta_y(k)
\left(
\hat{c}_{k,B}^\dagger \hat{c}_{k,A}
-
\mathrm{h.c.}
\right)
}_{H_{\mathrm{cross}}}
\end{equation}

Usando resultados conocidos, vemos que no existe corrección a primer
orden mientras que a segundo orden se encuentra una renormalización
de las bandas:

\begin{equation}
\Delta H^{(2)}
=
-
\sum_{k,\sigma}
\frac{
\left(g\Delta_z(k)\right)^2
}{
\hbar\omega_c
}
\hat{c}_{k,\sigma}^\dagger
\hat{c}_{k,\sigma}
\end{equation}

Con esto, logramos desacoplar efectivamente las bandas inferior y
superior en cada subespacio de igual número de fotones. Se observa
que este tratamiento es válido siempre que

\begin{equation}
\epsilon_{k,\sigma}
\gg
\frac{
\left(g\Delta_z(k)\right)^2
}{
\hbar\omega_c
}
\end{equation}

A partir de aquí nos enfocamos en el caso resonante donde la banda
de abajo desplazada un fotón coincide en energía con la banda de
arriba en ausencia de fotones, por lo cual:

\begin{equation}
\tilde{\epsilon}_{k,B}
-
\tilde{\epsilon}_{k,A}
\sim
\hbar\omega_c
\end{equation}

Pasando al picture de interacción y desechando los términos
contrarrotantes encontramos el siguiente hamiltoniano:

\begin{align}
\hat{H}_{RWA}
&=
\sum_{k,\sigma}
\tilde{\epsilon}_{k,\sigma}
\hat{c}_{k,\sigma}^\dagger
\hat{c}_{k,\sigma}
+
\hbar\omega_c \hat{a}^\dagger \hat{a}
\nonumber
\\
&\quad
+
i\sum_k
\Delta_y(k)
\left(
\hat{a}
\hat{c}_{k,B}^\dagger
\hat{c}_{k,A}
-
\hat{a}^\dagger
\hat{c}_{k,A}^\dagger
\hat{c}_{k,B}
\right)
\end{align}

donde efectivamente nos quedamos con la parte del hamiltoniano donde
cada emisión o absorción de un fotón lleva de una banda a la otra.

Definimos:

\begin{equation}
\hat{N}_\sigma
=
\sum_k
\hat{c}_{k,\sigma}^\dagger
\hat{c}_{k,\sigma}
\end{equation}

Se puede demostrar que:

\begin{equation}
\left[
\hat{H}_{RWA},
\hat{N}_A + \hat{N}_B
\right]
=
0
\end{equation}

\begin{equation}
\left[
\hat{H}_{RWA},
\hat{a}^\dagger \hat{a} + \hat{N}_B
\right]
=
0
\end{equation}

\begin{equation}
\left[
\hat{H}_{RWA},
\hat{a}^\dagger \hat{a} - \hat{N}_A
\right]
=
0
\end{equation}

y a su vez estos operadores conmutan entre sí, por lo cual sirven
para facilitar la diagonalización del problema.

\end{document}\documentclass[11pt]{article}

\usepackage[utf8]{inputenc}
\usepackage{amsmath,amssymb,bm}

\begin{document}

\section*{Hamiltoniano efectivo para una cadena SSH acoplada por un fotón}

Partimos del hamiltoniano

\begin{equation}
H = H_{SSH}(t_1, t_2) + g(a + a^\dagger) H_{SSH}(t_1, -t_2)
\end{equation}

con condiciones periódicas. Diagonalizamos la parte puramente material:

\begin{equation}
H_{SSH}(t_1, t_2)
=
\sum_k |\Delta_0(k)| \,
\psi_k^\dagger \hat{\sigma}_z^k \psi_k
=
\sum_{k,\sigma}
\epsilon_{k,\sigma}
\hat{c}_{k,\sigma}^\dagger
\hat{c}_{k,\sigma}
\end{equation}

con

\begin{equation}
\Delta_0(k) = t_1 + t_2 e^{-ik},
\end{equation}

\begin{equation}
\epsilon_{k,\sigma} = (-1)^\sigma |\Delta_0(k)|.
\end{equation}

Para esta diagonalización se hicieron dos pasos. Primero un cambio de gauge:

\begin{equation}
\hat{c}_{k,A}^\dagger \hat{c}_{k,B}
\rightarrow
e^{-i\varphi_0(k)}
\hat{c}_{k,A}^\dagger \hat{c}_{k,B}
\end{equation}

con

\begin{equation}
\varphi_0(k)
=
\operatorname{Arg}\left(\Delta_0(k)\right).
\end{equation}

Segundo, un cambio de variables:

\begin{equation}
\hat{c}_{k,A}
=
\frac{
\hat{d}_{k,A} + \hat{d}_{k,B}
}{\sqrt{2}}
\end{equation}

\begin{equation}
\hat{c}_{k,B}
=
\frac{
\hat{d}_{k,A} - \hat{d}_{k,B}
}{\sqrt{2}}
\end{equation}

En este cambio de variables las matrices de Pauli transforman a:

\begin{equation}
\sigma_x^k \rightarrow \sigma_z^k
\end{equation}

\begin{equation}
\sigma_y^k \rightarrow -\sigma_y^k
\end{equation}

\begin{equation}
\sigma_z^k \rightarrow \sigma_x^k
\end{equation}

Tras todo esto, el hamiltoniano de acoplamiento transformó a:

\begin{equation}
H_{SSH}(t_1,-t_2)
=
\sum_k
\psi_k^\dagger
\left(
\Delta_z(k)\hat{\sigma}_z^k
+
\Delta_y(k)\hat{\sigma}_y^k
\right)
\psi_k
\end{equation}

con

\begin{equation}
\Delta_z(k)
=
\frac{
t_1^2 - t_2^2
}{
|\Delta_0(k)|
}
\end{equation}

\begin{equation}
\Delta_y(k)
=
\frac{
2t_1 t_2 \sin(k)
}{
|\Delta_0(k)|
}.
\end{equation}

Se puede mostrar que el módulo de $\Delta_y(k)$ se anula en
$k = 0,\pi$, es una función impar y además tiene un máximo para
$\pi/2 < |k| < \pi$.

A continuación realizamos una transformación de Schrieffer--Wolff
para eliminar la parte diagonal en materia de este acoplamiento.
Definimos:

\begin{equation}
\hat{H} = \hat{H}_0 + \hat{V}
\end{equation}

con

\begin{equation}
\hat{V}
=
\sum_{k,\sigma}
\Delta_z^\sigma(k)
\hat{c}_{k,\sigma}^\dagger
\hat{c}_{k,\sigma},
\end{equation}

\begin{equation}
\Delta_z^\sigma(k)
=
(-1)^\sigma \Delta_z(k),
\end{equation}

y la parte que queremos conservar del hamiltoniano:

\begin{equation}
\hat{H}_0
=
H_{SSH}(t_1,t_2)
+
\hbar\omega_c \hat{a}^\dagger \hat{a}
+
g(a+a^\dagger)
\underbrace{
\sum_k
i\Delta_y(k)
\left(
\hat{c}_{k,B}^\dagger \hat{c}_{k,A}
-
\mathrm{h.c.}
\right)
}_{H_{\mathrm{cross}}}
\end{equation}

Usando resultados conocidos, vemos que no existe corrección a primer
orden mientras que a segundo orden se encuentra una renormalización
de las bandas:

\begin{equation}
\Delta H^{(2)}
=
-
\sum_{k,\sigma}
\frac{
\left(g\Delta_z(k)\right)^2
}{
\hbar\omega_c
}
\hat{c}_{k,\sigma}^\dagger
\hat{c}_{k,\sigma}
\end{equation}

Con esto, logramos desacoplar efectivamente las bandas inferior y
superior en cada subespacio de igual número de fotones. Se observa
que este tratamiento es válido siempre que

\begin{equation}
\epsilon_{k,\sigma}
\gg
\frac{
\left(g\Delta_z(k)\right)^2
}{
\hbar\omega_c
}
\end{equation}

A partir de aquí nos enfocamos en el caso resonante donde la banda
de abajo desplazada un fotón coincide en energía con la banda de
arriba en ausencia de fotones, por lo cual:

\begin{equation}
\tilde{\epsilon}_{k,B}
-
\tilde{\epsilon}_{k,A}
\sim
\hbar\omega_c
\end{equation}

Pasando al picture de interacción y desechando los términos
contrarrotantes encontramos el siguiente hamiltoniano:

\begin{align}
\hat{H}_{RWA}
&=
\sum_{k,\sigma}
\tilde{\epsilon}_{k,\sigma}
\hat{c}_{k,\sigma}^\dagger
\hat{c}_{k,\sigma}
+
\hbar\omega_c \hat{a}^\dagger \hat{a}
\nonumber
\\
&\quad
+
i\sum_k
\Delta_y(k)
\left(
\hat{a}
\hat{c}_{k,B}^\dagger
\hat{c}_{k,A}
-
\hat{a}^\dagger
\hat{c}_{k,A}^\dagger
\hat{c}_{k,B}
\right)
\end{align}

donde efectivamente nos quedamos con la parte del hamiltoniano donde
cada emisión o absorción de un fotón lleva de una banda a la otra.

Definimos:

\begin{equation}
\hat{N}_\sigma
=
\sum_k
\hat{c}_{k,\sigma}^\dagger
\hat{c}_{k,\sigma}
\end{equation}

Se puede demostrar que:

\begin{equation}
\left[
\hat{H}_{RWA},
\hat{N}_A + \hat{N}_B
\right]
=
0
\end{equation}

\begin{equation}
\left[
\hat{H}_{RWA},
\hat{a}^\dagger \hat{a} + \hat{N}_B
\right]
=
0
\end{equation}
++
\begin{equation}
\left[
\hat{H}_{RWA},
\hat{a}^\dagger \hat{a} - \hat{N}_A
\right]
=
0
\end{equation}

y a su vez estos operadores conmutan entre sí, por lo cual sirven
para facilitar la diagonalización del problema.

\end{document}